\documentclass[../../main.tex]{subfiles}% 注意这里的文件路径不能用 ./main.tex ,否则用latexmk编译子文件会报错
\graphicspath{{\subfix{./image/}}{\subfix{../../image/}}} % 指定图片引用路径,后续可以直接使用图片文件名
% 如果图片存放在上级目录,则使用 ../../image/ 作为备用路径

% 例如:
% \begin{figure}[H]
% \centering
% \includegraphics[scale=0.3]{图.png}
% \caption{}
% \label{figure:图}
% \end{figure}
% 注意:上述\label{}一定要放在\caption{}之后,否则引用图片序号会只会显示??.

\begin{document}

\section{曲面的存在性定理}

\begin{definition}
设$D \subset E^2$为$E^2$中的单连通区域，给定二次微分形式
\begin{align}
\varphi = g_{\alpha\beta}\mathrm{d}u^\alpha\mathrm{d}u^\beta,\quad \psi = b_{\alpha\beta}\mathrm{d}u^\alpha\mathrm{d}u^\beta,
\label{eq:::IO22JIOWAIJ2.1}
\end{align}
其中$g_{\alpha\beta}$在$D$上至少二阶连续可微，$b_{\alpha\beta}$至少一阶连续可微，满足$g_{\alpha\beta}=g_{\beta\alpha}$，$b_{\alpha\beta}=b_{\beta\alpha}$，且矩阵$(g_{\alpha\beta})$正定。记$(g^{\alpha\beta})$为$(g_{\alpha\beta})$的逆矩阵，定义：
\begin{align}
\Gamma_{\gamma\alpha\beta} &= \frac{1}{2}\left(\frac{\partial g_{\gamma\beta}}{\partial u^\alpha}+\frac{\partial g_{\alpha\gamma}}{\partial u^\beta}-\frac{\partial g_{\alpha\beta}}{\partial u^\gamma}\right),
\label{eq:::IO22JIOWAIJ2.2}\\
\Gamma^\gamma_{\alpha\beta} &= g^{\gamma\delta}\Gamma_{\delta\alpha\beta},
\label{eq:::IO22JIOWAIJ2.3}\\
R^\delta_{\alpha\beta\gamma} &= \frac{\partial \Gamma^\delta_{\alpha\beta}}{\partial u^\gamma}-\frac{\partial \Gamma^\delta_{\alpha\gamma}}{\partial u^\beta}+\Gamma^\eta_{\alpha\beta}\Gamma^\delta_{\eta\gamma}-\Gamma^\eta_{\alpha\gamma}\Gamma^\delta_{\eta\beta},
\label{eq:::IO22JIOWAIJ2.4}\\
R_{\alpha\delta\beta\gamma} &= g_{\delta\eta}R^\eta_{\alpha\beta\gamma}.
\label{eq:::IO22JIOWAIJ2.5}
\end{align}
\end{definition}

\begin{theorem}
若\eqref{eq:::IO22JIOWAIJ2.1}给出的二次微分形式$\varphi,\psi$满足\textbf{Gauss-Codazzi方程}
\begin{align}
\begin{split}
b_{11}b_{22}-(b_{12})^2 &= R_{1212},\\
\frac{\partial b_{11}}{\partial u^2}-\frac{\partial b_{12}}{\partial u^1} &= -b_{2\gamma}\Gamma^\gamma_{11}+b_{1\gamma}\Gamma^\gamma_{12},\\
\frac{\partial b_{21}}{\partial u^2}-\frac{\partial b_{22}}{\partial u^1} &= -b_{2\gamma}\Gamma^\gamma_{21}+b_{1\gamma}\Gamma^\gamma_{22},
\end{split}
\label{eq:::IO22JIOWAIJ2.6}
\end{align}
则对任意$(u_0^1,u_0^2) \in D$，存在邻域$U \subset D$，以及$E^3$中定义在$U$上的正则参数曲面$S: \boldsymbol{r}=\boldsymbol{r}(u^1,u^2)$，使得其第一基本形式、第二基本形式分别为$\varphi|_U,\psi|_U$；且任意两个满足条件的曲面可通过$E^3$的刚体运动重合。
\end{theorem}
\begin{remark}
定理的唯一性部分由前文曲面唯一性定理直接可得，故只需证明曲面的存在性；该定理揭示\textbf{Gauss-Codazzi方程}是给定两个二次微分形式为曲面基本形式的充分条件。
\end{remark}
\begin{remark}
将曲面视为一族标架的思想具有基础性意义，可将曲面存在性定理纳入标架族存在性定理的框架中研究。
\end{remark}
\begin{proof}

利用$\varphi,\psi$的系数及其导数，构造一阶线性齐次偏微分方程组
\begin{align}
\begin{cases}
\displaystyle\frac{\partial \boldsymbol{r}}{\partial u^\beta} = \boldsymbol{r}_\beta,\\[6pt]
\displaystyle\frac{\partial \boldsymbol{r}_\alpha}{\partial u^\beta} = \Gamma^\gamma_{\alpha\beta}\boldsymbol{r}_\gamma + b_{\alpha\beta}\boldsymbol{n},\\[6pt]
\displaystyle\frac{\partial \boldsymbol{n}}{\partial u^\beta} = -b^\gamma_\beta \boldsymbol{r}_\gamma,
\end{cases}
\label{eq:::IO22JIOWAIJ2.7}
\end{align}
未知向量函数为$\boldsymbol{r},\boldsymbol{r}_1,\boldsymbol{r}_2,\boldsymbol{n}$，共12个分量，自变量为$u^1,u^2$。

一阶偏微分方程组有解的充要条件为相容性条件
\begin{align*}
\begin{cases}
\displaystyle\frac{\partial}{\partial u^\gamma}\left(\frac{\partial \boldsymbol{r}}{\partial u^\beta}\right) = \frac{\partial}{\partial u^\beta}\left(\frac{\partial \boldsymbol{r}}{\partial u^\gamma}\right),\\[6pt]
\displaystyle\frac{\partial}{\partial u^\gamma}\left(\frac{\partial \boldsymbol{r}_\alpha}{\partial u^\beta}\right) = \frac{\partial}{\partial u^\beta}\left(\frac{\partial \boldsymbol{r}_\alpha}{\partial u^\gamma}\right),\\[6pt]
\displaystyle\frac{\partial}{\partial u^\gamma}\left(\frac{\partial \boldsymbol{n}}{\partial u^\beta}\right) = \frac{\partial}{\partial u^\beta}\left(\frac{\partial \boldsymbol{n}}{\partial u^\gamma}\right).
\end{cases}
\end{align*}
该相容性条件等价于Gauss-Codazzi方程\eqref{eq:::IO22JIOWAIJ2.6}，故方程组\eqref{eq:::IO22JIOWAIJ2.7}完全可积。

对任意$(u_0^1,u_0^2) \in D$，取初值$\boldsymbol{r}^0,\boldsymbol{r}_1^0,\boldsymbol{r}_2^0,\boldsymbol{n}^0$，存在邻域$U \subset D$与解
\begin{align}
\begin{split}
\boldsymbol{r} &= \boldsymbol{r}(u^1,u^2),\\
\boldsymbol{r}_1 &= \boldsymbol{r}_1(u^1,u^2),\\
\boldsymbol{r}_2 &= \boldsymbol{r}_2(u^1,u^2),\\
\boldsymbol{n} &= \boldsymbol{n}(u^1,u^2),
\end{split}
\label{eq:::IO22JIOWAIJ2.8}
\end{align}
满足初值条件
\begin{align}
\begin{cases}
\boldsymbol{r}(u_0^1,u_0^2) = \boldsymbol{r}^0,\\
\boldsymbol{r}_1(u_0^1,u_0^2) = \boldsymbol{r}_1^0,\\
\boldsymbol{r}_2(u_0^1,u_0^2) = \boldsymbol{r}_2^0,\\
\boldsymbol{n}(u_0^1,u_0^2) = \boldsymbol{n}^0.
\end{cases}
\label{eq:::IO22JIOWAIJ2.9}
\end{align}

限定初值满足自然标架条件
\begin{align}
\begin{split}
\boldsymbol{r}_\alpha^0 \cdot \boldsymbol{r}_\beta^0 &= g_{\alpha\beta}(u_0^1,u_0^2),\\
\boldsymbol{r}_\alpha^0 \cdot \boldsymbol{n}^0 &= 0,\\
\boldsymbol{n}^0 \cdot \boldsymbol{n}^0 &= 1,\\
(\boldsymbol{r}_1^0,\boldsymbol{r}_2^0,\boldsymbol{n}^0) &> 0,
\end{split}
\label{eq:::IO22JIOWAIJ2.10}
\end{align}
以保证$\boldsymbol{r}_1,\boldsymbol{r}_2,\boldsymbol{n}$线性无关。

定义函数
\begin{align}
\begin{split}
f_{\alpha\beta}(u^1,u^2) &= \boldsymbol{r}_\alpha(u^1,u^2) \cdot \boldsymbol{r}_\beta(u^1,u^2) - g_{\alpha\beta}(u^1,u^2),\\
f_\alpha(u^1,u^2) &= \boldsymbol{r}_\alpha(u^1,u^2) \cdot \boldsymbol{n}(u^1,u^2),\\
f(u^1,u^2) &= \boldsymbol{n}(u^1,u^2) \cdot \boldsymbol{n}(u^1,u^2) - 1,
\end{split}
\label{eq:::IO22JIOWAIJ2.11}
\end{align}
由\eqref{eq:::IO22JIOWAIJ2.9}\eqref{eq:::IO22JIOWAIJ2.10}得初值
\begin{align}
f_{\alpha\beta}(u_0^1,u_0^2)=f_\alpha(u_0^1,u_0^2)=f(u_0^1,u_0^2)=0.
\label{eq:::IO22JIOWAIJ2.12}
\end{align}

将$f_{\alpha\beta},f_\alpha,f$求导，代入\eqref{eq:::IO22JIOWAIJ2.7}可得
\begin{align}
\begin{cases}
\displaystyle\frac{\partial f_{\alpha\beta}}{\partial u^\gamma} = \Gamma^\delta_{\gamma\alpha}f_{\delta\beta}+\Gamma^\delta_{\gamma\beta}f_{\delta\alpha}+b_{\gamma\alpha}f_\beta+b_{\gamma\beta}f_\alpha,\\[6pt]
\displaystyle\frac{\partial f_\alpha}{\partial u^\gamma} = -b^\delta_\gamma f_{\delta\alpha}+\Gamma^\delta_{\gamma\alpha}f_\delta+b_{\gamma\alpha}f,\\[6pt]
\displaystyle\frac{\partial f}{\partial u^\gamma} = -2b^\delta_\gamma f_\delta.
\end{cases}
\label{eq:::IO22JIOWAIJ2.13}
\end{align}

该方程组与初值条件和前文曲面唯一性定理的对应形式完全一致，故仅有零解，即
$$f_{\alpha\beta}(u^1,u^2)=f_\alpha(u^1,u^2)=f(u^1,u^2) \equiv 0,$$
因此在$U$上恒成立
\begin{align}
\begin{split}
\boldsymbol{r}_\alpha(u^1,u^2) \cdot \boldsymbol{r}_\beta(u^1,u^2) &= g_{\alpha\beta}(u^1,u^2),\\
\boldsymbol{r}_\alpha(u^1,u^2) \cdot \boldsymbol{n}(u^1,u^2) &= 0,\\
\boldsymbol{n}(u^1,u^2) \cdot \boldsymbol{n}(u^1,u^2) &= 1.
\end{split}
\label{eq:::IO22JIOWAIJ2.14}
\end{align}

计算混合积
\begin{align*}
(\boldsymbol{r}_1,\boldsymbol{r}_2,\boldsymbol{n})^2 &= \det\left(\begin{pmatrix}\boldsymbol{r}_1\\\boldsymbol{r}_2\\\boldsymbol{n}\end{pmatrix} \cdot (\boldsymbol{r}_1,\boldsymbol{r}_2,\boldsymbol{n})\right)\\
&= \det\begin{pmatrix}g_{11}&g_{12}&0\\g_{21}&g_{22}&0\\0&0&1\end{pmatrix}\\
&= \det(g_{\alpha\beta}) > 0,
\end{align*}
故$\boldsymbol{r}_1,\boldsymbol{r}_2,\boldsymbol{n}$处处线性无关，结合连续性得
\begin{align}
(\boldsymbol{r}_1(u^1,u^2),\boldsymbol{r}_2(u^1,u^2),\boldsymbol{n}(u^1,u^2)) > 0,
\label{eq:::IO22JIOWAIJ2.15}
\end{align}
即$\{\boldsymbol{r}_1,\boldsymbol{r}_2,\boldsymbol{n}\}$构成右手系，且
\begin{align}
\boldsymbol{n} = \frac{\boldsymbol{r}_1 \times \boldsymbol{r}_2}{|\boldsymbol{r}_1 \times \boldsymbol{r}_2|}.
\label{eq:::IO22JIOWAIJ2.16}
\end{align}

由\eqref{eq:::IO22JIOWAIJ2.7}第一式，$\boldsymbol{r}_1,\boldsymbol{r}_2$为曲面$\boldsymbol{r}=\boldsymbol{r}(u^1,u^2)$的切向量，线性无关故曲面正则；$\boldsymbol{n}$为单位法向量，结合\eqref{eq:::IO22JIOWAIJ2.14}与\eqref{eq:::IO22JIOWAIJ2.7}第二式，曲面的第一、第二基本形式恰好为$\varphi,\psi$.

\end{proof}

\begin{example}
验证: 由\eqref{eq:::IO22JIOWAIJ2.11}式定义的函数$f_{\alpha\beta}(u^1,u^2)$，$f_\alpha(u^1,u^2)$，$f(u^1,u^2)$满足一阶线性齐次偏微分方程组\eqref{eq:::IO22JIOWAIJ2.13}.
\end{example}

\begin{example}
判断下面给出的二次微分形式$\varphi,\psi$能否作为空间$E^3$中一块曲面的第一基本形式和第二基本形式? 说明理由.
\begin{enumerate}[(1)]
\item $\varphi=(\mathrm{d}u)^2+(\mathrm{d}v)^2,\quad \psi=(\mathrm{d}u)^2-(\mathrm{d}v)^2;$
\item $\varphi=(\mathrm{d}u)^2+\cos^2 u\,(\mathrm{d}v)^2,\quad \psi=\cos^2 u\,(\mathrm{d}u)^2+(\mathrm{d}v)^2.$
\end{enumerate}
\end{example}

\begin{example}
求曲面$S$的参数方程，使得它的第一基本形式和第二基本形式分别为
\begin{align*}
\mathrm{I}=(1+u^2)(\mathrm{d}u)^2+u^2\,(\mathrm{d}v)^2,\quad \mathrm{II}=\frac{1}{\sqrt{1+u^2}}\big((\mathrm{d}u)^2+u^2\,(\mathrm{d}v)^2\big).
\end{align*}
\end{example}

\begin{example}
已知二次微分形式$\varphi,\psi$如下:
\begin{align*}
\varphi=E(u,v)(\mathrm{d}u)^2+G(u,v)(\mathrm{d}v)^2,\quad \psi=\lambda(u,v)\cdot\varphi,
\end{align*}
其中函数$E(u,v)>0$，$G(u,v)>0$。若$\varphi,\psi$能够作为空间$E^3$中一块曲面的第一基本形式和第二基本形式，则函数$E(u,v)$，$G(u,v)$，$\lambda(u,v)$应该满足什么条件?

假定$E(u,v)=G(u,v)$，写出满足上述条件的函数$E(u,v)$，$G(u,v)$，$\lambda(u,v)$的具体表达式.
\end{example}



\end{document}